\subsection{Construções em \textbf{Trees}}
\label{subsec:trees-constructions}

\subsubsection{Produto}

\begin{proposition}
O produto $T_1 \times T_2$ de duas árvores binárias pode ser construído como a árvore cujos nós são os pares $(u, v)$ com $u \in T_1$ e $v \in T_2$, com raiz $(\text{root}_1, \text{root}_2)$, e onde $(u,v)$ tem filho esquerdo/direito se e só se $u$ e $v$ têm filho esquerdo/direito em simultâneo.
\end{proposition}

\begin{remark}
Esta construção é válida mas pode colapsar rapidamente para a árvore vazia se as duas árvores têm profundidades diferentes --- o produto nesta categoria é \emph{não trivial} e pode ser a árvore com um só nó (se as árvores têm alturas diferentes). Na prática, o produto existe mas pode ser menor do que qualquer dos fatores.
\end{remark}

\begin{lstlisting}[style=matlab, caption={Produto de duas árvores binárias (representação recursiva)}]
function T = tree_product(T1, T2)
    % Constroi o produto T1 x T2 recursivamente
    % Retorna [] se um dos nos nao tem correspondente no outro
    if isempty(T1) || isempty(T2)
        T = []; return;
    end
    T.value = [T1.value, T2.value];  % par de valores
    T.left  = tree_product(T1.left,  T2.left);
    T.right = tree_product(T1.right, T2.right);
end
\end{lstlisting}

\subsubsection{Coproduto}

\begin{proposition}
O coproduto (soma disjunta) $T_1 \sqcup T_2$ é a árvore obtida criando uma nova raiz com $T_1$ como sub-árvore esquerda e $T_2$ como sub-árvore direita.
\end{proposition}

\begin{lstlisting}[style=matlab, caption={Coproduto de duas árvores binárias}]
function T = tree_coproduct(T1, T2)
    T.value = [];     % nova raiz sem valor proprio
    T.left  = T1;
    T.right = T2;
end
\end{lstlisting}

\subsubsection{Pullback}

Dado um diagrama $T_1 \xrightarrow{f} T_3 \xleftarrow{g} T_2$, o pullback é a subárvore de $T_1 \times T_2$ formada pelos pares $(u,v)$ tais que $f(u) = g(v)$.

\begin{lstlisting}[style=matlab, caption={Pullback de duas árvores sobre uma terceira}]
function [PB, pi1, pi2] = tree_pullback(T1, T2, T3, f, g)
    % f: morfismo T1->T3 (vetor), g: morfismo T2->T3 (vetor)
    % Constroi os nos do pullback como pares (i,j) com f(i)==g(j)
    nodes = [];
    for i = 1:T1.n
        for j = 1:T2.n
            if f(i) == g(j)
                nodes(end+1, :) = [i, j];
            end
        end
    end
    PB.nodes = nodes;  % cada linha e um par (i,j)
    pi1 = nodes(:,1)'; % projecao para T1
    pi2 = nodes(:,2)'; % projecao para T2
    % A estrutura de arvore do pullback herda-se de T1 x T2
end
\end{lstlisting}

\subsubsection{Pushout}

O pushout de $T_1 \xleftarrow{f} T_0 \xrightarrow{g} T_2$ é o coproduto $T_1 \sqcup T_2$ quocientado pela relação que identifica $f(u)$ com $g(u)$ para cada nó $u \in T_0$. Em árvores binárias, isto corresponde a ``colar'' as duas árvores ao longo da imagem comum.

\subsubsection{Equalizador}

O equalizador de dois morfismos $f, g: T_1 \rightrightarrows T_2$ é a subárvore de $T_1$ induzida pelo conjunto de nós $E = \{u \in T_1 : f(u) = g(u)\}$, desde que $E$ forme uma subárvore (fechada para progenitores e para a raiz).

\begin{lstlisting}[style=matlab, caption={Equalizador de dois morfismos de árvores}]
function [E, inc] = tree_equalizer(T1, f, g)
    % E: subconjunto dos nos de T1 onde f e g coincidem
    mask = (f == g);
    E    = find(mask);    % indices dos nos do equalizador
    inc  = E;             % inclusao em T1
    % Nota: verificar que E e fechado para progenitores
end
\end{lstlisting}

\subsubsection{Coequalizador}

O coequalizador de $f, g: T_1 \rightrightarrows T_2$ é o quociente de $T_2$ pela menor relação de equivalência que identifica $f(u)$ com $g(u)$ para todo $u \in T_1$, desde que o quociente admita estrutura de árvore binária.

\begin{remark}
Nem sempre o quociente é uma árvore binária (pode violar a estrutura de ramificação). Nestes casos, o coequalizador não existe em \textbf{Trees}, o que ilustra que esta categoria não é cocompleta em geral.
\end{remark}